\begin{table*}[t]
\centering
\caption{Treatment-level performance results.}
\label{tab:treatment_results}
\renewcommand{\arraystretch}{1.15}
\resizebox{\textwidth}{!}{%
\begin{tabular}{p{2.8cm}cccp{1.6cm}p{3.4cm}p{2.4cm}}
\toprule
\textbf{Treatment} &
\textbf{Macro-P} &
\textbf{Macro-R} &
\textbf{Macro-F1} &
\textbf{Avg.\ pred.\ issues} &
\textbf{Primary characteristic} &
\textbf{Role} \\
\midrule

Single-LLM Family Mean
& 0.0769
& ---
& 0.0960
& ---
& Balanced single-model behaviour
& Reference \\

Equal-Weight Aggregation
& 0.0771
& 0.4077
& \textbf{0.1153}
& 144.8
& Highest recall; very large issue list
& Aggregation baseline \\

Roundtable No Weighting
& ---
& ---
& 0.0968
& 136.7
& Consensus without reliability-aware weighting
& Ablation \\

Full Method
& \textbf{0.0901}
& 0.1354
& 0.0711
& \textbf{15.9}
& Highest precision; strongest issue-list compression
& Proposed \\

\bottomrule
\end{tabular}%
}

\vspace{0.4em}
\begin{minipage}{0.97\textwidth}
\footnotesize
\textit{Overall comparison.} Friedman test on document-level Macro-F1 finds no statistically significant overall difference among the four treatments ($p>0.05$). Although the proposed Full Method does not achieve the highest Macro-F1, it produces substantially more compact issue lists while improving precision, indicating a selective issue-generation strategy rather than a high-recall detector.
\textit{Note.} Macro-P = Macro-Precision; Macro-R = Macro-Recall. All performance values are averaged over the 10 shared requirements documents. Cells marked ``---'' denote values not emphasized in the current paper and can be added if desired.
\end{minipage}
\end{table*}
